% ============================================================================
% Fichier  : front/01_avant_propos.tex — Avant-propos
% Version  : 1.0 -- novembre 2026 (proposition à relire par l'auteur)
% ============================================================================

\chapter*{Avant-propos}
\addcontentsline{toc}{chapter}{Avant-propos}
\markboth{Avant-propos}{Avant-propos}

Une banque peut être certifiée et compromise la même année. Ce paradoxe
apparent, qui ouvre ce livre, n'a rien d'exceptionnel~: il naît chaque fois
qu'une organisation confond ce qu'elle a fait vérifier avec ce qu'elle est
réellement, et ce qu'elle a regardé avec ce qui existe. Les normes, les
méthodes de gestion des risques et l'audit sont les outils dont on dispose
pour réduire cet écart~; ils peuvent aussi le masquer, quand on les emploie
pour rassurer plutôt que pour savoir.

Ce cours s'adresse aux étudiants du master professionnel en management des
systèmes numériques, spécialité sécurité des systèmes numériques, de la
Chaire UNESCO d'Afrique centrale, et aux praticiens qui exercent ou
exerceront les métiers de responsable de la sécurité des systèmes
d'information, d'auditeur en cybersécurité, de gestionnaire des risques
cyber, de responsable GRC ou de consultant en mise en œuvre de systèmes de
management. Il suppose les bases de la cybersécurité~; il ne suppose aucune
connaissance préalable des normes, de la gestion des risques ni de l'audit.

Il a trois partis pris. Le premier est de relier, plutôt que de juxtaposer,
les normes, le risque et l'audit~: une même organisation les vit ensemble,
et les échecs naissent souvent entre eux. Le deuxième est d'ancrer chaque
notion dans le contexte camerounais et africain~: le droit de la
cybersécurité, la régulation bancaire de la CEMAC, les organismes publics
et leurs contrôleurs, au côté des normes internationales. Le troisième est
de mesurer ce que l'on ne sait pas~: chaque cote, chaque conclusion, chaque
chiffre du livre dit sur quoi il repose, et l'instrument \OUI{}~/ \NON{}~/
\NEK{}, issu du programme de recherche de l'auditologie, donne à cette
exigence une forme mesurable.

Les trois organisations qui accompagnent le lecteur, la BANQUE SANAGA,
LiZbiZ et l'ANPP, sont fictives, comme leurs dirigeants, leurs prestataires,
les dates et les montants. Les mécanismes, les normes et les textes cités
sont exacts, avec, pour chaque texte, le degré de vérification indiqué à
l'annexe~\ref{ann:textes}. Les normes sous licence sont paraphrasées, jamais
citées.

Le livre est accompagné d'un livret des corrigés, qui explique et annote
chaque exercice, et de NDAP SUKULU, sur lequel les fiches d'auto-évaluation
et les ateliers du cours permettent de pratiquer chaque chapitre.

\bigskip
\hfill\begin{minipage}{0.5\textwidth}
\raggedleft
\IfFileExists{../../Logos/signature_avant_propos.png}{\includegraphics[height=1.6cm]{../../Logos/signature_avant_propos.png}}{\textit{MINKA MI NGUIDJOÏ Thierry Emmanuel}}\\
Yaoundé, novembre 2026
\end{minipage}
